\medskip

\emph{Les parties {\rm A} et {\rm B} peuvent être traitées indépendamment}

\medskip

\textbf{Partie A}

\medskip

Le plan  est ramené à un repère orthogonal.

On a représenté ci-dessous la courbe d'une fonction $f$ définie et deux fois dérivable sur $\R$, ainsi que celle de sa dérivée $f'$ et de sa dérivée seconde $f''$.

\begin{center}
\psset{xunit=1.5cm,yunit=1.0cm,algebraic=true,dimen=middle,dotstyle=o,dotsize=5pt 0,linewidth=2.pt,arrowsize=2pt 3}
\begin{pspicture*}(-0.5,-4.05)(8,5.2)
\psgrid[gridlabels=0pt,subgriddiv=1,gridwidth=0.2pt](0,-4)(8,5)
\psaxes[labelFontSize=\scriptstyle,xAxis=true,yAxis=true,,ticksize=-2pt 0,subticks=2]{->}(0,0)(0,-3.95)(8,5)
\psline[linewidth=4.pt](-3.320081944682136,4.772338854767012)(-1.5600159659972166,4.772338854767012)
\psplot[linewidth=1.25pt,linecolor=red,plotpoints=2000]{0}{8}{4.0/(1.0+2.718281828^(-3.0*x+12.0))}\uput[r](5,3.6){\red $\mathcal{C}_2$}
\psplot[linewidth=1.25pt,linecolor=blue,plotpoints=2000]{0}{8}{4.0*3.0*2.71828^(-3.0*x+12.0)/(2.0*2.71828^(-3.0*x+12.0)+(2.71828^(-3.0*x+12.0))^(2.0)+1.0)}\uput[r](5,0.6){\blue $\mathcal{C}_3$}
\psplot[linewidth=1.25pt,plotpoints=2000]{0}{8}{(-4.0*3.0^(2.0)*2.71828^(-3.0*x+12.0)+4.0*3.0^(2.0)*(2.71828^(-3.0*x+12.0))^(2.0))/(3.0*2.71828^(-3.0*x+12.0)+(2.71828^(-3.0*x+12.0))^(3.0)+3.0*(2.71828^(-3.0*x+12.0))^(2.0)+1.0)}\uput[r](5,-1.6){$\mathcal{C}_1$}
\end{pspicture*}
\end{center}

\begin{enumerate}
\item Déterminer, en justifiant votre choix, quelle courbe correspond à quelle fonction.
\item Déterminer, avec la précision permise par le graphique, le coefficient directeur de la tangente à la courbe $\mathcal{C}_2$ au point d'abscisse 4.
\item Donner avec la précision permise par le graphique, l'abscisse de chaque point d'inflexion de la courbe $\mathcal{C}_1$.
\end{enumerate}

\medskip

\textbf{Partie B}

\medskip

Soit un réel $k$ strictement positif.

On considère la fonction $g$ définie sur $\R$ par:
\[g(x) = \dfrac{4}{1 + \text{e}^{- kx}}.\]

\smallskip

\begin{enumerate}
\item  Déterminer les limites de $g$ en $+ \infty$ et en $-\infty$,
\item  Prouver que $g'(0) = k$.
\item  En admettant le résultat ci-dessous obtenu avec un logiciel de calcul formel,
prouver que la courbe de $g$ admet un point d'inflexion au point d'abscisse 0.
\begin{center}
\begin{tabular}{|c l|}\hline
\multicolumn{2}{|l|}{$\triangleright$ Calcul formel}\\ \hline
&g(x) =4/(1 + e{}\^{}(-kx))\\
1&\\
&$\to g(x) = \dfrac{4}{\text{e}^{-kx} + 1}$\\ \hline
&Simplifier$(g''(x))$\\
2&\\
&$\to g''(x) = - 4\text{e}^{kx}\left(\text{e}^{kx} - 1\right)\dfrac{k^2}{\left(\text{e}^{kx} + 1 \right)^3}$\\ \hline
\end{tabular}
\end{center}
\end{enumerate}

\bigskip
